\item La función $C(t) = K(e^{-at} - e^{-bt})$, en donde $a$, $b$ y $K$ son constantes positivas, y $b > a$, se emplea para representar la concentración en el tiempo $t$ de una droga o fármaco inyectado en la sangre. (a) Demuestre que $\lim_{t \to \infty} C(t) = 0$. (b) Encuentre $C'(t)$, la rapidez a la que el producto se elimina de la circulación. (c) Determine cuándo dicha velocidad es igual a 0.
